%%%%%%%%%%%%
%
% $Autor: Wings $
% $Datum: 2019-03-05 08:03:15Z $
% $Pfad: Maintenance.tex $
% $Version: 4250 $
% !TeX spellcheck = en_GB/de_DE
% !TeX encoding = utf8
% !TeX root = manual 
% !TeX TXS-program:bibliography = txs:///biber
%
%%%%%%%%%%%%

\chapter{Wartung}

\section{Einmal jeden Monat}

\begin{itemize}
	
	\item \textbf{Vorspannung kontrollieren} \\
	Vorspannung des Transportbandes prüfen. Bei erhöhtem Spiel die Vorspannung des Transportbandes erhöhen.
	
	\item \textbf{Verschleißprüfung} \\
	Sichtprüfung des Transportbands auf Schnitte oder Risse. Wenn Beschädigungen vorhanden sind, ist das Transportband zu erneuern.
	
	\item \textbf{Geräuschprüfung} \\
	Laufruhe des Transportbands sowie Geräuschentwicklung der Komponenten testen. Hierfür das elektrische Tischförderband bei unterschiedlichen Geschwindigkeiten betreiben. Die Geräuschentwicklung sollte mit erhöhter Geschwindigkeit gleichmäßig zunehmen. Bei erhöhter Lautstärke Vorspannung prüfen und bei Bedarf verringern. Wenn das Problem weiterhin auftritt, Lagerung und Antrieb soweit möglich kontrollieren.
	
	\item \textbf{Prüfung der Spurhaltigkeit} \\
	Zum Prüfen das Elektrische Tischförderband mit Maximalgeschwindigkeit und ohne Last aktivieren. Dann einseitig die Einstellmuttern der Vorspannungsverstellung nutzen, um das Transportband mittig auf den Rollen auszurichten.
	
	\item \textbf{Prüfung der Sensorik} \\
	Während aktivem Transportbandlauf ohne Last den Sensor durch einen beliebigen Gegenstand aktivieren. Die Reaktion des Förderbandes muss sofortig erfolgen. Wenn dies nicht der Fall ist, Neustart durchführen. Bei bleibendem Fehler Software und Elektronik prüfen lassen.
	
\end{itemize}

\section{Halbjährlich}

\begin{itemize}
	
	\item \textbf{Prüfung der Laufruhe} \\
	Prüfung des Bandes bei 50 \% und später 100 \% Geschwindigkeit auf Laufruhe. Wenn unruhiger Lauf oder Schwingungen im Betrieb festgestellt werden, sind die Transportbandvorspannung, die Laufrollen sowie die Lagerung zu kontrollieren. Wenn diese trotz unrundem Lauf fehlerfrei sind, Transportband erneuern.
	
	\item \textbf{Last- und Stabilitätsprüfung} \\
	Förderband bei maximaler Geschwindigkeit und maximaler Last mehrfach prüfen. Geladene Objekte müssen stabil und ruckfrei transportiert werden. Die Geräuschentwicklung sollte sich auch nach mehreren Durchgängen nicht signifikant erhöhen. Wenn dies nicht der Fall ist, Vorspannung verringern und Komponenten auf Beschädigungen prüfen.
	
\end{itemize}

\section{Einstellen und Austauschen von Komponenten}

\begin{itemize}
	
	\item \textbf{Vorspannung ändern} \\
	Zur Änderung der Vorspannung die Einstellmuttern an den Seiten des Förderbandes drehen, bis die gewünschte Spannung erreicht ist. Achtung: Bei zu hoher Vorspannung kann es zu Beschädigungen der Antriebskomponenten kommen. Ist die Vorspannung zu niedrig kann das Transportband durchrutschen oder von den Führungen gleiten.
	
	\item \textbf{Transportband wechseln} \\
	Um das Transportband zu wechseln die Schnellspannvorrichtung lösen und das Transportband seitlich entnehmen. Vor Einbau des Neuteils die Vorspanneinrichtung etwas lösen, dann das neue Transportband auflegen und die Schnellspannvorrichtung wieder schließen. Anschließend Vorspannung wie beschrieben einstellen.
	
\end{itemize}
